% Generated draft: requires mathematical and source review.
\documentclass[11pt]{article}
\usepackage[margin=25mm]{geometry}
\usepackage{fontspec}
\setmainfont{Latin Modern Roman}
\usepackage{amsmath,amssymb,mathtools,unicode-math}
\setmathfont{Latin Modern Math}
\usepackage{enumitem,needspace,xurl,hyperref}
\setlength{\parindent}{0pt}
\setlength{\parskip}{5pt}
\setlength{\emergencystretch}{3em}
\newcommand{\N}{\mathbb{N}}
\newcommand{\Z}{\mathbb{Z}}
\newcommand{\Q}{\mathbb{Q}}
\newcommand{\R}{\mathbb{R}}
\newcommand{\C}{\mathbb{C}}
\newcommand{\F}{\mathbb{F}}
\newcommand{\A}{\mathbb{A}}
\newcommand{\PP}{\mathbb{P}}
\newcommand{\E}{\mathbb{E}}
\newcommand{\eps}{\varepsilon}
\newcommand{\dd}{\,\mathrm{d}}
\begin{document}
\section*{UnsolvedMath ID: 30002510}
\textbf{Source problem code:} catalogue entry
\subsection{Definitions and mathematical statement}
The supplied problem asks to construct ``good spines of Teichmüller space''; its normalized formulation asks for geometrically natural, mapping-class-group-equivariant spines. Neither ``good'' nor ``geometrically natural'' is defined in the record, and the surface type and range of genera are not specified.

For a conditional, standard closed-surface interpretation, fix an integer \(g\geq 2\) and a closed, connected, oriented surface \(S_g\) of genus \(g\). Let \(\mathcal T_g\) be the Teichmüller space of marked compact Riemann surfaces \((X,f)\), where \(f:S_g\to X\) is an orientation-preserving homeomorphism. Two markings \((X,f)\) and \((Y,k)\) represent the same point if there is a conformal isomorphism \(a:X\to Y\) such that \(a\circ f\) is isotopic to \(k\). Let \(\operatorname{Mod}(S_g)\) be the group of isotopy classes of orientation-preserving self-homeomorphisms of \(S_g\), acting by \([h]\cdot[X,f]=[X,f\circ h^{-1}]\).

One precise possible meaning of an equivariant spine is an invariant subset \(K_g\subseteq\mathcal T_g\) admitting a continuous equivariant strong deformation retraction
\[
H_g:\mathcal T_g\times[0,1]\longrightarrow\mathcal T_g.
\]
Explicitly, for every \(x\in\mathcal T_g\), \(k\in K_g\), \(t\in[0,1]\), and \(\gamma\in\operatorname{Mod}(S_g)\), require
\[
H_g(x,0)=x,\qquad H_g(x,1)\in K_g,\qquad H_g(k,t)=k,
\]
\[
\gamma K_g=K_g,\qquad H_g(\gamma x,t)=\gamma H_g(x,t).
\]
Under this interpretation, the construction problem is to produce such spines by a geometrically natural construction for the intended genera. This is only a conditional formalization: the source does not specify whether an equivariant strong deformation retraction is intended, what additional structure a spine must have, or what further properties distinguish a ``good'' or ``geometrically natural'' spine. Thus the displayed conditions specify a possible topological requirement but do not fully determine the source's intended target.
\subsection{Short English statement}
Construct geometrically natural spines of Teichmüller space compatible with the mapping class group action. A standard closed-surface interpretation and explicit equivariant deformation-retraction conditions are given above, but the intended surface range and the additional meaning of ``good'' remain unspecified.
\subsection{Sources}
\begin{itemize}
\item UnsolvedMath local registry, permanent ID \(30002510\), title: \emph{Good Spines of Teichmüller Space}; exact source problem code: \texttt{catalogue entry}. The record gives proposed year \(2014\) and names no proposer. Its original wording is ``Construct good spines of Teichmuller space T\_g.'' The supplied normalized wording adds ``geometrically natural'' and ``mapping-class-group-equivariant.''
\item Supplied report citation: \emph{New Trends in Teichmüller Theory and Mapping Class Groups} (2014), Oberwolfach Reports, pp.\ 395--452, DOI \texttt{10.4171/owr/2014/07}. \url{https://doi.org/10.4171/owr/2014/07}. The background describes these reports as documenting mathematical workshops and their open-problem discussions.
\item Supplied literature source, labeled ``Oberwolfach Report 07/2014'': \url{https://doi.org/10.4171/owr/2014/07}.
\item Supplied literature source, labeled ``Thurston spine dimension estimates'': \url{https://arxiv.org/abs/2310.15618}.
\item Literature assessment dated exactly \texttt{2026-08-21}: the supplied curation labels the problem open and states that several explicit spines and dimension estimates are available, while a satisfactory general construction of good spines remains open. This is supplied dated literature triage derived from the OWR curation, not an independently verified assessment; time-sensitive status and cited claims require independent verification.
\item The registry labels the statement \texttt{corrected\_verified} and describes its changes as normalization of notation, typography, and wording without changing mathematical content. These are supplied registry assertions, not independent verification of the formulation or its equivalence to the report.
\end{itemize}
\end{document}
